\newpage
\lecture{8}{24 March.}{}
\begin{observation}
    Let \(G = (V, E)\) and suppose \(G\) has a hamilton cycle. Then, for every \(\varnothing \neq S \subseteq V\), 
    \[
        \underbrace{c(G - S)}_{\# \text{ of cpts}} \le \vert S \vert
    \]   
    Let $C$ be a Hamiltonian cycle of $G$. Consider removing the vertices in $S$ from $C$. Since $C$ is a cycle passing through all vertices of $G$, deleting $S$ breaks $C$ into a collection of vertex-disjoint paths, each of which lies entirely in $G - S$.

Each such path is contained in a connected component of $G - S$. Moreover, every component of $G - S$ must contain at least one of these path segments, since $C$ visits every vertex of $G$.

Now fix a component $H$ of $G - S$. Along the cycle $C$, the vertices of $H$ appear consecutively in one or more segments. Each maximal segment of $C$ contained in $H$ must be entered from $S$ and exited back to $S$. Therefore, for each such segment, there are at least two edges of $C$ between $S$ and $H$ (one entering and one leaving).

In particular, each component $H$ of $G - S$ is incident with at least two edges of $C$ that go between $S$ and $V(H)$. Hence, the total number of edges of $C$ between $S$ and $G - S$ is at least
\[
    2 \, c(G - S).
\]

On the other hand, each vertex in $S$ has degree exactly $2$ in the cycle $C$. Therefore, the total number of edges of $C$ incident to vertices in $S$ is at most
\[
    2|S|.
\]
All edges between $S$ and $G - S$ are among these edges, so
\[
    \#\{\text{edges between } S \text{ and } G - S \text{ in } C\} \le 2|S|.
\]

Combining the two inequalities, we obtain
\[
    2\,c(G - S) \le 2|S|,
\]
which implies
\[
    c(G - S) \le |S|.
\]
\end{observation}

\begin{definition}[toughness]
    Given \(t \in \mathbb{R} \), we say a graph \(G\) is \(t\)-tough if 
    \[
        \vert S \vert \ge t c(G - S) 
    \]   
    for all \(S \subseteq V(G)\) and \(S \neq \varnothing \).  
\end{definition}

Hence, the observation above states that the Hamiltonicity implies \(1\)-tough. 

\begin{conjecture}[Chvatal, 1973]
    There is some \(t\) s.t. \(t\)-tough implies Hamiltonicity.  
\end{conjecture}

Construction of Bauer-Broersma-Veldman (2000): families of non-Hamiltonian graphs with toughness \(\left( \frac{9}{4} - o(1) \right) \). 

\begin{note}
    Toughness of \(C_n \coloneqq 1\). 
\end{note}

\section{Complexity}

\begin{definition}
    A decision problem is a problem of answer yes/no.
\end{definition}

\begin{remark}
More formally, a decision problem can be viewed as a language $L \subseteq \{0,1\}^*$, where an input $x$ is a \textbf{yes-instance} if $x \in L$, and a \textbf{no-instance} otherwise.
\end{remark}

\begin{definition}[polynomial time]
    A decision problem with an algorithm to find the answer in polynomial time in the size of the input means there exists a constant $\alpha$ s.t. if input has size $n$, then we can find the answer in time $O(n^{\alpha})$.   
\end{definition}

\begin{remark}
The class of all decision problems solvable in polynomial time is denoted by $P$. 
Polynomial time is considered ``efficient'' or ``tractable'' in theoretical computer science.
\end{remark}

\begin{definition}[Nondeterministic polynomial time]
    "Yes" answers come with a certificate that can be verified in polynomial time.
\end{definition}

\begin{remark}
More precisely, a decision problem is in $NP$ if there exists a polynomial-time algorithm $V(x,c)$ (called a \textbf{verifier}) such that:
\begin{itemize}
    \item if $x$ is a yes-instance, then there exists a certificate $c$ with $|c| \le \text{poly}(|x|)$ such that $V(x,c)=1$,
    \item if $x$ is a no-instance, then for all $c$, $V(x,c)=0$.
\end{itemize}
\end{remark}

\begin{eg}
    In Hamiltonicity problem, the certificate is the hamilton cycle. 
\end{eg}

\begin{remark}
Given a sequence of vertices, we can check in polynomial time whether it forms a Hamiltonian cycle by verifying:
\begin{itemize}
    \item all vertices appear exactly once,
    \item consecutive vertices are adjacent,
    \item the first and last vertices are also adjacent.
\end{itemize}
\end{remark}

\begin{claim}
    $P \subseteq NP$. 
\end{claim}
\begin{explanation}
    To check the certificate, ignore it and solve the problem from scratch.
\end{explanation}

\begin{remark}
Indeed, if a problem can be solved in polynomial time, then we can verify it in polynomial time by simply recomputing the answer. Hence every problem in $P$ is also in $NP$.
\end{remark}

\begin{definition}[co-NP]
    co-NP problem is the class of decision problems where No instances have a certificate that can be verified in polynomial time.
\end{definition}

\begin{remark}
Equivalently, a problem $L$ is in co-$NP$ if its complement $\overline{L}$ is in $NP$.
\end{remark}

\begin{eg}[co-NP]
\hfill
    \begin{itemize}
        \item anything in $P$. 
        \item Is this graph not Hamiltonian?
        \item Does a graph $G$ has a perfect matching? (This problem is in $NP \cap \text{co-}NP$ since for yes instances the certificate is the perfect matching itself and for no instances we can use Tutte's theorem, i.e. the certificate is $S \subseteq V$ s.t. $o(G - S) > |S|$).  
    \end{itemize}
\end{eg}

\begin{remark}
For many problems (such as Hamiltonicity), it is unknown whether they belong to co-$NP$. 
In particular, we do not know whether Hamiltonicity has efficiently verifiable certificates for \emph{non}-existence.
\end{remark}

\begin{question}
    Does $P = NP$? 
\end{question}

Common belief: $P \neq NP$.

\begin{remark}
This is one of the most important open problems in computer science. It asks whether every efficiently verifiable problem is also efficiently solvable.
\end{remark}

\begin{definition}[NP-hard]
    A problem that any NP-problem can be reduced to in polynomial time is called NP-hard.
\end{definition}

\begin{remark}
More precisely, a problem $H$ is NP-hard if for every problem $L \in NP$, there exists a polynomial-time reduction $L \le_p H$.
Note that NP-hard problems are not required to be in $NP$ (they may even be undecidable).
\end{remark}

\begin{definition}[NP-complete]
    $NP$-complete $= NP$-hard $\cap NP$.   
\end{definition}

\begin{remark}
NP-complete problems are the ``hardest'' problems in $NP$. If any NP-complete problem can be solved in polynomial time, then $P = NP$.
\end{remark}

Karp (1972) provided an initial list of many NP-complete problem, including
\begin{itemize}
    \item Hamiltonicity. 
    \item 3-Colorability. 
    \item 3-SAT.
\end{itemize}

\begin{remark}
These problems are polynomially equivalent under reductions, meaning that an efficient algorithm for any one of them would yield efficient algorithms for all NP problems.
\end{remark}

\section{Maximum matching problem}
\begin{question}
    Given a graph \(G\), what is the size of its largest matching? 
\end{question}

This is an optimization problem (OPT). 

The related decision problem is: Is the size of the largest matching \(\ge k\)?

\begin{theorem}[Hall, min-max version]
    Let \(G = (X \cup Y, E)\) be a bipartite graph. Then, 
    \[
        \max \vert M \vert = \min _{S \subseteq X} \vert X \vert - \vert S \vert + \vert N(S) \vert,
    \] 
    where \(\vert M \vert \) is the number of edges in a matching. 
\end{theorem}

\begin{proof}
We first show
\[
    \max |M| \le \min_{S \subseteq X} \big( |X| - |S| + |N(S)| \big).
\]

Fix any $S \subseteq X$, and let $M$ be any matching in $G$.

Each vertex in $N(S)$ can be matched to at most one vertex in $S$, so at most $|N(S)|$ vertices of $S$ are saturated by $M$. Hence at least
\[
    |S| - |N(S)|
\]
vertices of $S$ are unsaturated.

Therefore, the total number of vertices in $X$ that are saturated by $M$ is at most
\[
    |X| - (|S| - |N(S)|) = |X| - |S| + |N(S)|.
\]

Since this holds for every $S \subseteq X$, we obtain
\[
    |M| \le \min_{S \subseteq X} \big( |X| - |S| + |N(S)| \big).
\]

Taking maximum over all matchings $M$, we get
\[
    \max |M| \le \min_{S \subseteq X} \big( |X| - |S| + |N(S)| \big).
\]

\medskip

Now we prove the reverse inequality.

Define the \emph{deficiency} of a set $S \subseteq X$ by
\[
    |S| - |N(S)|.
\]
Let
\[
    t := \max_{S \subseteq X} |S| - |N(S)| = \max_{S \subseteq X} \big( |S| - |N(S)| \big).
\]

Then for every $S \subseteq X$, we have
\[
    |S| - |N(S)| \le t,
\]
which implies
\[
    |N(S)| \ge |S| - t.
\]

We now construct a new bipartite graph $G' = (X \cup (Y \cup Y'), E')$ by adding a set $Y'$ of $t$ new vertices, each adjacent to every vertex in $X$.

For any nonempty $S \subseteq X$, we have
\[
    N_{G'}(S) = N_G(S) \cup Y',
\]
and hence
\[
    |N_{G'}(S)| = |N_G(S)| + t \ge (|S| - t) + t = |S|.
\]

Thus Hall's condition holds in $G'$, and by Hall's theorem there exists a matching in $G'$ that saturates all vertices in $X$.

Since $|Y'| = t$, at most $t$ vertices of $X$ are matched to vertices in $Y'$. Therefore, at least $|X| - t$ vertices of $X$ are matched to vertices in $Y$.

Removing all edges incident to $Y'$, we obtain a matching in $G$ of size at least $|X| - t$. Hence
\[
    \max |M| \ge |X| - t.
\]

Finally, observe that
\[
    |X| - t = |X| - \max_{S} (|S| - |N(S)|)
    = \min_{S} \big( |X| - |S| + |N(S)| \big).
\]

Therefore,
\[
    \max |M| \ge \min_{S \subseteq X} \big( |X| - |S| + |N(S)| \big).
\]

\medskip

Combining the two inequalities, we conclude
\[
    \max |M| = \min_{S \subseteq X} \big( |X| - |S| + |N(S)| \big).
\]
\end{proof}

\begin{remark}
    This shows that "Does a bipartite graph have a matching of size \(k\)?" is in NP and co-NP.
\end{remark}

\begin{definition}
    A vertex cover of a graph is a set \(C \subseteq V\) of vertices such that for every edge \(e \in E\), \(e \cap C \neq \varnothing\).   
\end{definition}

\begin{definition}
    Given a graph \(G\), we write 
    \begin{align*}
        \alpha ^{\prime} (G) &= \max \left\{ \vert M \vert : M \subseteq E \text{ is a matching}   \right\} \\
        \beta (G) &= \min \left\{ \vert C \vert: C \subseteq V \text{ is a vertex cover}   \right\}. 
    \end{align*}
\end{definition}

\begin{claim}
    For any \(G\), \(\alpha ^{\prime} (G) \le \beta (G)\), i.e. for any matching \(M\) and vertex cover \(C\), \(\vert M \vert \le \vert C \vert  \).     
\end{claim}

\begin{explanation}
    Let $M$ be any matching in $G$, and let $C$ be any vertex cover of $G$.

By definition, every edge in $M$ must be incident to at least one vertex in $C$, since $C$ is a vertex cover.

On the other hand, since $M$ is a matching, no two edges in $M$ share a common endpoint. Therefore, each vertex in $C$ can cover at most one edge of $M$.

Hence, to cover all edges in $M$, we need at least $|M|$ vertices in $C$. That is,
\[
    |C| \ge |M|.
\]

Since this holds for all matchings $M$ and all vertex covers $C$, we conclude
\[
    \alpha'(G) \le \beta(G).
\]
\end{explanation}

\begin{remark}
    We don't necessarily have \(\alpha ^{\prime} (G) = \beta (G)\) (consider \(K_3\)). 
\end{remark}

\begin{theorem}[Konig]
    If \(G=(X \cup  Y, E)\) is bipartite, then 
    \[
        \alpha ^{\prime} (G) = \beta (G).
    \] 
\end{theorem}

\begin{proof}
    We already know \(\alpha ^{\prime} (G) \le \beta (G)\). Thus, we need to show \(\alpha ^{\prime} (G) \ge \beta (G)\). Let \(C \subseteq V\) be a smallest vertex cover, then we need to find a matching of size \(\vert C \vert \). We will use Hall's theorem to find a matching in \((C \cap X) \cup (Y \setminus C)\) saturating \(C \cap X\) and a matching in \((C \cap Y) \cup (X \setminus C)\) saturating \(C \cap Y\). Then, the union of these matchings is a matching of size 
    \[
        \vert C \cap X \vert + \vert C \cap Y \vert = \vert C \vert.   
    \] 
    If not, then WLOG there is no matching in \((X \cap C) \cup (Y \setminus C)\) saturating \(X \cap C\). By Hall's theorem, \(\exists S \subseteq X \cap C\) s.t. 
    \[
        \vert N_{Y \setminus C}(S) \vert < \vert S \vert,  
    \]          
    but then 
    \[
        C^{\prime} = (C \setminus S) \cup N_{Y \setminus C}(S)
    \]
    is a vertex cover of smaller size, which is a contradiction.

    \begin{center}
\begin{tikzpicture}[x=1cm,y=1cm,>=Latex, every node/.style={font=\small}]
  \node[font=\bfseries\large] at (-4.8,3.25) {$X$};
  \node[font=\bfseries\large] at (4.8,3.25) {$Y$};

  \draw[thick] (-4.8,0.25) ellipse (1.55 and 2.35);
  \draw[thick] (4.8,0.25) ellipse (1.55 and 2.35);

  \node[circle,draw,fill=red!20,minimum size=7mm] (x1) at (-5.2,1.9) {$x_1$};
  \node[circle,draw,fill=red!20,minimum size=7mm] (x2) at (-4.4,1.1) {$x_2$};
  \node[circle,draw,fill=orange!35,minimum size=7mm] (x3) at (-5.2,-0.2) {$x_3$};
  \node[circle,draw,fill=orange!35,minimum size=7mm] (x4) at (-4.4,-1.0) {$x_4$};
  \node at (-4.8,0.25) {$\vdots$};

  \node[circle,draw,fill=red!20,minimum size=7mm] (y1) at (4.4,1.9) {$y_1$};
  \node[circle,draw,fill=red!20,minimum size=7mm] (y2) at (5.2,1.1) {$y_2$};
  \node[circle,draw,fill=blue!25,minimum size=7mm] (y3) at (4.4,-0.2) {$y_3$};
  \node[circle,draw,fill=orange!35,minimum size=7mm] (y4) at (5.2,-1.0) {$y_4$};
  \node[circle,draw,fill=orange!35,minimum size=7mm] (y5) at (4.8,-1.8) {$y_5$};
  \node at (4.8,0.25) {$\vdots$};

  \draw[very thick] (x1) -- (y3);
  \draw[very thick] (x2) -- (y3);

  \draw[gray!70] (x3) -- (y1);
  \draw[gray!70] (x3) -- (y4);
  \draw[gray!70] (x4) -- (y2);
  \draw[gray!70] (x4) -- (y5);

  \draw[dashed,rounded corners] (-5.75,0.55) rectangle (-3.9,2.35);
  \node[anchor=east] at (-5.8,2.45) {$S\subseteq X\cap C$};

  \draw[dashed,rounded corners] (3.95,-0.65) rectangle (4.85,0.25);
  \node[anchor=west] at (4.95,0.45) {$N_{Y\setminus C}(S)$};

  \node[red!70!black] at (-4.8,2.65) {$X\cap C$};
  \node[red!70!black] at (4.8,2.65) {$Y\cap C$};
  \node[orange!60!black] at (-4.8,-2.35) {$X\setminus C$};
  \node[orange!60!black] at (4.8,-2.35) {$Y\setminus C$};

  \node[align=center] at (0,-2.10) {$|N_{Y\setminus C}(S)|<|S|$};
  \node[align=center] at (0,-2.55) {$C'=(C\setminus S)\cup N_{Y\setminus C}(S)$};
  \node[align=center] at (0,-3.00) {$|C'|<|C|\Rightarrow$ contradiction};

  \node[draw,rounded corners,align=left,anchor=north west] at (-7.0,-2.65) {
    \textcolor{red!70!black}{紅色：在 $C$ 中}\\
    \textcolor{blue!60!black}{藍色：$N_{Y\setminus C}(S)$}\\
    \textcolor{orange!60!black}{橘色：不在 $C$ 中}
  };
\end{tikzpicture}
\end{center}
\end{proof}


